\chapter{Hipergrafos, \(\alpha\)-aciclicidad y \emph{join trees}}
\label{cap:hipergrafos}

\begin{intuicion}
Hasta ahora vimos joins binarios (\(R \bowtie S\)) y joins
multi-way ``planos''. Para entender Yannakakis necesitamos una
herramienta más: el \emph{hipergrafo de la consulta}. Dado un join
de \(n\) relaciones, su hipergrafo tiene los atributos como nodos y
cada relación como una hiperarista (subset de atributos). La
propiedad mágica es la \(\alpha\)-aciclicidad: si el hipergrafo de
la consulta es \(\alpha\)-acíclico, existe un \emph{join tree} sobre
las relaciones, y Yannakakis resuelve el join en tiempo lineal en
input + output. Todas las cadenas, estrellas, árboles puros son
\(\alpha\)-acíclicos. Triángulo, \(K_4\) y otros cíclicos no lo son.
Este capítulo da las definiciones, muestra el algoritmo GYO de
detección, y construye los join trees que el matcher de Yannakakis
recorre.
\end{intuicion}

\section{Hipergrafos: del join al \emph{query graph}}

\begin{definition}[Hipergrafo de una consulta]
Dado un join \(Q = R_1 \bowtie R_2 \bowtie \dots \bowtie R_m\)
sobre los atributos \(\mathcal{A} = \bigcup_l \mathrm{attrs}(R_l)\),
su \emph{hipergrafo} es \(H(Q) = (\mathcal{A}, F)\) donde \(F =
\{F_1, F_2, \dots, F_m\}\) y \(F_l = \mathrm{attrs}(R_l)\) (cada
relación es una hiperarista).
\end{definition}

\begin{ejemplo}[Hipergrafo de una cadena de 3]
\(Q = R(A, B) \bowtie S(B, C) \bowtie T(C, D)\).

Hipergrafo: \(\mathcal{A} = \{A, B, C, D\}\), \(F = \{F_R = \{A,
B\}, F_S = \{B, C\}, F_T = \{C, D\}\}\).

\begin{figure}[H]
\centering
\begin{tikzpicture}[
    every node/.style={font=\small},
    attr/.style={draw=accent, thick, circle, minimum size=8mm,
                 fill=accentlight, inner sep=1pt},
    edge/.style={draw=accent, thick, dashed}
]
\node[attr] (A) at (0,0) {A};
\node[attr] (B) at (2,0) {B};
\node[attr] (C) at (4,0) {C};
\node[attr] (D) at (6,0) {D};

\draw[edge, rounded corners] ($(A.north west) + (-0.3, 0.3)$) rectangle ($(B.south east) + (0.3, -0.3)$);
\node[font=\scriptsize, accent] at (1, 0.6) {$F_R$};

\draw[edge, rounded corners] ($(B.north west) + (-0.3, -0.5)$) rectangle ($(C.south east) + (0.3, -1.1)$);
\node[font=\scriptsize, accent] at (3, -1.4) {$F_S$};

\draw[edge, rounded corners] ($(C.north west) + (-0.3, 0.3)$) rectangle ($(D.south east) + (0.3, -0.3)$);
\node[font=\scriptsize, accent] at (5, 0.6) {$F_T$};
\end{tikzpicture}
\caption{Hipergrafo de la cadena de 3 relaciones. Cada hiperarista
(rectángulo dashed) es una relación. Los atributos compartidos
(\(B\) entre \(F_R\) y \(F_S\); \(C\) entre \(F_S\) y \(F_T\))
son los join keys.}
\label{fig:hipergrafo-cadena}
\end{figure}
\end{ejemplo}

\begin{ejemplo}[Hipergrafo del triángulo]
\(Q = R(A, B) \bowtie S(B, C) \bowtie T(A, C)\).

Hipergrafo: \(\mathcal{A} = \{A, B, C\}\), \(F = \{F_R = \{A, B\},
F_S = \{B, C\}, F_T = \{A, C\}\}\).

Cada par de hiperaristas comparte exactamente un atributo:
\(F_R \cap F_S = \{B\}\), \(F_S \cap F_T = \{C\}\),
\(F_R \cap F_T = \{A\}\). Eso forma un \emph{ciclo} en el sentido
hipergráfico: \(F_R \to F_S\) (vía \(B\)), \(F_S \to F_T\) (vía
\(C\)), \(F_T \to F_R\) (vía \(A\)).
\end{ejemplo}

\section{\(\alpha\)-aciclicidad}

La definición clásica es operacional, vía el algoritmo GYO.

\subsection{El algoritmo GYO}

\begin{definition}[Reducción GYO]
Dado un hipergrafo \(H = (\mathcal{A}, F)\), aplicar repetidamente
las dos reglas:

\begin{enumerate}
    \item \textbf{Eliminar atributo aislado}: si un atributo
    \(a \in \mathcal{A}\) aparece en exactamente una hiperarista
    \(F_l\), eliminarlo de \(F_l\) y de \(\mathcal{A}\).
    \item \textbf{Eliminar hiperarista contenida}: si una hiperarista
    \(F_l\) está contenida en otra \(F_{l'}\) (\(F_l \subseteq
    F_{l'}\)), eliminar \(F_l\).
\end{enumerate}

Si después de aplicar las reglas hasta no poder más, \(F = \emptyset\),
el hipergrafo es \emph{\(\alpha\)-acíclico}. Si no, es cíclico.
\end{definition}

\begin{ejemplo}[GYO sobre la cadena]
\(F_R = \{A, B\}\), \(F_S = \{B, C\}\), \(F_T = \{C, D\}\).

\begin{itemize}
    \item Atributo \(A\) aparece sólo en \(F_R\). Eliminarlo:
    \(F_R = \{B\}\).
    \item Atributo \(D\) aparece sólo en \(F_T\). Eliminarlo:
    \(F_T = \{C\}\).
    \item Ahora \(F_R = \{B\} \subseteq F_S = \{B, C\}\). Eliminar
    \(F_R\).
    \item Ahora \(F_T = \{C\} \subseteq F_S = \{B, C\}\). Eliminar
    \(F_T\).
    \item Quedó \(F = \{F_S\}\). Atributos \(B\) y \(C\) aparecen
    sólo en \(F_S\); eliminarlos. \(F_S = \emptyset\); eliminar
    \(F_S\).
    \item \(F = \emptyset\). \textbf{Cadena es \(\alpha\)-acíclica.}
\end{itemize}
\end{ejemplo}

\begin{ejemplo}[GYO sobre el triángulo]
\(F_R = \{A, B\}\), \(F_S = \{B, C\}\), \(F_T = \{A, C\}\).

\begin{itemize}
    \item Atributo \(A\): aparece en \(F_R\) y \(F_T\). No
    aislado.
    \item \(B\): aparece en \(F_R\) y \(F_S\). No aislado.
    \item \(C\): aparece en \(F_S\) y \(F_T\). No aislado.
    \item No hay hiperarista contenida en otra (todas tienen 2
    atributos distintos).
    \item Ninguna regla aplica. \(F \neq \emptyset\). \textbf{Triángulo
    es \(\alpha\)-cíclico.}
\end{itemize}
\end{ejemplo}

\begin{nota}
\(\alpha\)-aciclicidad NO coincide con la noción intuitiva de
``ciclo en el grafo''. Hay hipergrafos que ``parecen tener ciclos''
pero son \(\alpha\)-acíclicos en el sentido GYO. Y al revés.
La definición es operacional y vale lo que vale para Yannakakis:
si GYO reduce a vacío, Yannakakis funciona en lineal.
\end{nota}

\section{El \emph{join tree}}

\begin{theorem}[Join tree iff \(\alpha\)-aciclicidad]
Un hipergrafo \(H\) es \(\alpha\)-acíclico \emph{si y sólo si}
existe un \emph{join tree}: un árbol \(T\) cuyos nodos son las
hiperaristas \(F_l\), tal que para cualquier atributo \(a \in
\mathcal{A}\), el conjunto de nodos \(\{F_l : a \in F_l\}\) forma
un sub-árbol conexo de \(T\).
\end{theorem}

\begin{proof}[Idea]
Forward: si \(H\) es \(\alpha\)-acíclico, GYO da un orden de
eliminación; el join tree se construye conectando cada hiperarista
eliminada a la que la ``contuvo'' en su eliminación. Backward: dado
un join tree, GYO siempre puede reducir empezando por las hojas.
\(\square\)
\end{proof}

\begin{ejemplo}[Join tree de la cadena]
Para la cadena \(R(A, B) \bowtie S(B, C) \bowtie T(C, D)\), el
join tree es:
\[
F_R \;\;-\!-\!-\;\; F_S \;\;-\!-\!-\;\; F_T
\]
\noindent
\(B\) aparece en \(F_R\) y \(F_S\): conexo (camino directo). \(C\)
aparece en \(F_S\) y \(F_T\): conexo. \(A\) aparece sólo en
\(F_R\); \(D\) sólo en \(F_T\). Ambos sub-árboles triviales.
\end{ejemplo}

\begin{ejemplo}[Join tree de una estrella]
\(Q = R_0(A, B_1) \bowtie R_1(A, B_2) \bowtie R_2(A, B_3)\) (tres
relaciones que comparten el atributo central \(A\)).

Join tree:
\begin{figure}[H]
\centering
\begin{tikzpicture}[
    every node/.style={draw=accent, thick, rectangle, rounded corners=1mm,
                       minimum size=7mm, fill=accentlight, font=\small}
]
\node (R0) at (0, 1) {$F_{R_0}$};
\node (R1) at (-2, -0.5) {$F_{R_1}$};
\node (R2) at (2, -0.5) {$F_{R_2}$};

\draw[accent, thick] (R0) -- (R1);
\draw[accent, thick] (R0) -- (R2);
\end{tikzpicture}
\caption{Join tree de una estrella de tres relaciones. \(A\) aparece
en todas: sub-árbol = todo el árbol (conexo).}
\end{figure}
\end{ejemplo}

\section{Por qué \(\alpha\)-aciclicidad importa}

\begin{theorem}[Yannakakis 1981]
Si \(H(Q)\) es \(\alpha\)-acíclico, el join \(Q\) puede computarse
en tiempo \(O(|Q_{\mathrm{in}}| + |Q_{\mathrm{out}}|)\), donde
\(|Q_{\mathrm{in}}| = \sum_l |R_l|\) y \(|Q_{\mathrm{out}}| = |Q|\).
\end{theorem}

\noindent
Lineal en input + output. Compará con el AGM bound
\(O(m^{\rhostar})\): para \(\alpha\)-acíclicas con
\(\rhostar > 1\), Yannakakis es estrictamente mejor (no en peor caso
porque output puede ser grande, pero el \emph{trabajo
constructivo} es lineal).

El truco es la \emph{semi-join reducer} en dos pasadas sobre el
join tree (capítulo~\ref{cap:yannakakis-chain}): elimina del input
todas las tuplas que no participan de ningún match completo,
\emph{antes} de empezar a enumerar. Después la enumeración es
trivial.

\section{La pieza estructural en GraphHunter}

\begin{itemize}
    \item Hipótesis \emph{shape chain}: hipergrafo es cadena.
    \(\alpha\)-acíclico. Yannakakis chain dispatcha
    (capítulo~\ref{cap:yannakakis-chain}).
    \item Hipótesis \emph{shape graph} con estructura de árbol /
    estrella: hipergrafo \(\alpha\)-acíclico. Yannakakis tree-walk
    dispatcha (capítulo~\ref{cap:yannakakis-tree}).
    \item Hipótesis \emph{shape graph} con ciclos (\(K_3\), \(K_4\)):
    cíclico. LFTJ dispatcha si es \(K_4\)
    (capítulo~\ref{cap:lftj}); si no, fallback al DFS.
\end{itemize}

\noindent
La detección de \(\alpha\)-aciclicidad ocurre en el \emph{planner}
(capítulo~\ref{cap:planner}): aplica GYO al hipergrafo de la
hipótesis. \(O(n_{\mathrm{atrs}} + n_{\mathrm{hipearistas}})\) que
para hipótesis típicas (\(\leq 6\) pasos) es trivial.

\section{Construcción concreta del join tree}

Para shape graph en GraphHunter, el planner construye el join tree
con un BFS sobre el hipergrafo a partir de una hiperarista
\emph{root} (típicamente la del primer paso de la hipótesis):

\begin{lstlisting}[language=Rust]
pub struct JoinTree {
    pub nodes: Vec<usize>,         // step indices (orden de visita)
    pub parent: Vec<Option<usize>>, // parent[i] = padre en el árbol
    pub shared_var: Vec<VarId>,    // variable compartida con el padre
}

impl JoinTree {
    pub fn build(hyp: &Hypothesis) -> Self {
        let hg = build_hypergraph(hyp);
        let mut visited = HashSet::new();
        let mut nodes = Vec::new();
        let mut parent = Vec::new();
        let mut shared_var = Vec::new();

        let root = 0;  // first step
        let mut queue = VecDeque::from([(root, None, None)]);
        while let Some((node, p, sv)) = queue.pop_front() {
            if visited.insert(node) {
                nodes.push(node);
                parent.push(p);
                shared_var.push(sv.unwrap_or(VarId::None));
                for (next, var) in hg.adjacent(node) {
                    if !visited.contains(&next) {
                        queue.push_back((next, Some(nodes.len() - 1), Some(var)));
                    }
                }
            }
        }

        Self { nodes, parent, shared_var }
    }
}
\end{lstlisting}

\noindent
El \texttt{enumerate\_tree} (capítulo~\ref{cap:yannakakis-tree})
recorre estos nodos en orden y, en cada uno, decide qué \emph{bag}
consultar según qué variables ya están bound.

\section{Resumen}

\begin{itemize}
    \item Hipergrafo de una consulta: atributos como nodos,
    relaciones como hiperaristas.
    \item \(\alpha\)-aciclicidad: el algoritmo GYO reduce a vacío.
    Operacionalmente, ``no hay ciclos en el sentido hipergráfico''.
    \item Cadena, estrella, árboles: \(\alpha\)-acíclicos. Triángulo,
    \(K_4\): cíclicos.
    \item Join tree: existe sii \(\alpha\)-aciclicidad. Conexidad
    de sub-árbol por atributo.
    \item Yannakakis (capítulo siguiente) corre en \(O(\)input +
    output\()\) sobre \(\alpha\)-acíclicos.
    \item GraphHunter dispatcha Yannakakis chain o tree-walk para
    hipótesis \(\alpha\)-acíclicas (chain o graph respectivamente);
    LFTJ para \(K_4\) cíclico; DFS para todo lo demás.
\end{itemize}

\begin{ejercicios}
\begin{enumerate}
    \item \starred{} Aplicá GYO al hipergrafo de la hipótesis
    \texttt{c -[Auth]-> h1; c -[Auth]-> h2; c -[Auth]-> h3}
    (estrella de tres). ¿Es \(\alpha\)-acíclico?

    \item Aplicá GYO al hipergrafo de \(K_4\) (6 hiperaristas
    binarias). Mostrá que es cíclico.

    \item \starred{} Construí un hipergrafo \(\alpha\)-acíclico
    de 4 relaciones. ¿Cuántos posibles join trees existen?

    \item ¿Por qué la noción de \(\alpha\)-aciclicidad NO
    coincide con la de ``no hay ciclos en el grafo''? Construí un
    contraejemplo: hipergrafo con ``ciclos visuales'' que es
    \(\alpha\)-acíclico, y otro al revés.

    \item ¿Cómo cambia el join tree de una hipótesis chain si la
    hipótesis es ``mismo vértice repetido'' (cadena
    \(p_0 \to p_1 \to p_2 \to p_0\))? ¿Sigue siendo \(\alpha\)-acíclico?

    \item Si una hipótesis tiene 5 pasos en chain pura, ¿cuántos
    nodos tiene el join tree? ¿Cuál es la altura del árbol?
\end{enumerate}
\end{ejercicios}
